\documentclass[a4paper,12pt]{report}

\usepackage[%
  % romanian / russian
  russian,%
  % an: курсовая
  % licenta: дипломная
  % licenta_practica_1/2/3: практика в 1/2/3 году
  % master: магистерская
  % master_practica_2: практика на 2 году мастерата
  master% 
]{config}

% design_jocurilor
% informatica
% informatica_aplicata
% informatica_stiinte_educatiei
\specialty{informatica_aplicata}

\newcommand{\uniGroupName}{MIA2201}

\newcommand{\authorName}{Nume Prenume}
\newcommand{\conducatorName}{Nume Prenume}
\newcommand{\conducatorTitleRo}{asistent universitar} % doctor, conferențiar universitar

\newcommand{\authorNameRu}{Фамилия Имя Отчество}
\newcommand{\conducatorNameRu}{Фамилия Имя}
\newcommand{\conducatorTitleRu}{ассистент университета} % доктор, преподаватель университета

\newcommand{\docTitleEng}{Thesis Name} % Нужно только на мастерате
\newcommand{\docTitleRo}{Explorarea formatului PNG}
\newcommand{\docTitleRu}{Исследование формата PNG}

\newcommand{\yearOfWriting}{\the\year}
\newcommand{\conferencesList}{Студенческая конференция номер НОМЕР, \yearOfWriting~года}
\newcommand{\github}{\url{https://github.com/USER/REPO}}
\newcommand{\outputDate}{\today}

% none - не использовал ИИ
% text - использовал исключительно для техноредактирования (правописание, организация материала, форматирование)
% assisted - использовал для еще чего-то
\newcommand{\aiUsage}{assisted}
% Пояснение использования ИИ, только если выбрали "assisted", на русском или на румынском.
% Если "assisted" не выбрано, это будет проигнорировано!
\newcommand{\aiAssistanceExplanation}{Использовал для}

% Добавьте сокращения здесь.
% Добавляйте только упомянутые сокращения.
% В дальнейшем всегда пишите \ac{PNG} вместо просто PNG.
\acro{PNG}{Portable Network Graphics}
\acro{PR}{Pull Request}
\acro[url_different_name]{URL}{Uniform Resource Locator}
\acro{LONG TEST}{VERY LONG NAME VERY LONG NAME VERY LONG NAME VERY LONG NAME VERY LONG NAME VERY LONG NAME VERY LONG NAME VERY LONG NAME}

\begin{document}

\titlePage

\clearpage
\tableofcontents

% Только для магистратуры!
\annotationChapter{ru}

\textbf{к \docTypeRuDatelny{} \enquote{\docTitleRu{}}, студента \authorNameRu{}, группа \uniGroupName{}, программа обучения \programulDeStudiiRu.}

\textbf{Структура диссертации.}
Диссертация состоит из: Введения, \chapterCount{} глав, Общих выводов и рекомендаций, Библиографии из \bibliographyEntryCount{} наименований.
Основной текст занимает \usefulPageCount{} страниц и включает \anexeCount{} приложения.

\textbf{Ключевые слова:}
% \acrshort делает так, чтобы PNG не раскрывалось в ее полное название на данном этапе,
% но оставалось кликабельным.
\textit{\acrshort{PNG}, слово 1, слово 2}

\textbf{Актуальность.}

Где используются технологии, почему тема важна.

\textbf{Цель и задачи исследования.}

Цель = что вы пытаетесь доказать в этой диссертации.

Задачи = шаги, через которые вы это доказываете.

\textbf{Ожидаемые и полученные результаты.}
Их можно резюмировать в:
(1) изучение чего-то
(2) проектирование чего-то
(3) реализация чего-то.

\textbf{Важная решённая проблема.}
Важная решённая проблема исследования состоит в ...

\textbf{Значимость, новизна и оригинальность.}
Работа отличается ...

\textbf{Практическая ценность, внедрение результатов и публикация.}
В результате было получено ...

Весь исходный код проекта доступен на GitHub по следующей ссылке: \github.

Полученные результаты были представлены на \textbf{\conferencesList}\cite{self}.

% Только для магистратуры!
\annotationChapter{ro}

\textbf{la \docTypeRoNedefinit{} \enquote{\docTitleRo{}}, a studentului \authorName{}, grupa \uniGroupName{}, programul de studii \programulDeStudiiRo.}

\textbf{Structura tezei.}
Teza constă din: Introducere, \chapterCount{} capitole, Concluzii generale și recomandări, Bibliografie \bibliographyEntryCount{} titluri.
Textul de bază cuprinde \usefulPageCount{} de pagini și \anexeCount{} anexe.

\textbf{Cuvinte-cheie:}
\textit{\acrshort{PNG}, cuvânt 1, cuvânt 2}

\textbf{Actualitatea.}

Unde se folosesc tehnologiile, de ce e importantă temă.

\textbf{Scopul și obiectivele cercetării.}

Scopul = ce încercați să demonstrați prin această teză.

Obiectivele = lucrurile prin care demonstrați asta.

\textbf{Rezultatele preconizate și obținute.}
Acestea rezumă în:
(1) studierea ceva
(2) proiectarea ceva
(3) implementarea ceva.

\textbf{Problema importantă rezolvată.}
Problema importantă rezolvată în cercetare constă în ...

\textbf{Semnificația, noutatea și originalitatea.}
Lucrarea se deosebește prin ...

\textbf{Valoarea aplicativă, implementarea rezultatelor și publicarea.}
Ca rezultat s-a obținut ...

Toate codurile sursă a proiectului pot fi accesate pe GitHub după următorul link: \github.

Rezultatele obținute au fost raportate la \textbf{\conferencesList}\cite{self}.

% Только для магистратуры!
\annotationChapter{en}

\textbf{of the \docTypeEng{} \enquote{\docTitleEng{}} of the student \authorName{}, group \uniGroupName{}, \programulDeStudiiEng{} study program.}

\textbf{The structure of the thesis.}
The thesis is made up of: Introduction, \chapterCount{} chapters,
General conclusions and recommendations, Bibliography of \bibliographyEntryCount{} titles.
Base text takes up \usefulPageCount{} pages and has \anexeCount{} appendicies.

\textbf{Keywords:}
\textit{}

\textbf{Relevance.}

\textbf{Purpose and objectives of this research.}

\textbf{The planned and obtained results.}

\textbf{Important resolved problem.}

\textbf{Significance, novelty, and originality.}

\textbf{Applicative value, implementation of results, and publication.}

The entire source code can be obtained by accessing the GitHub repository
by the following \gls{url_different_name}: \github. 

The obtained results were reported at \textbf{\conferencesList}\cite{self}.

\acronymsChapter

\introChapter

\textbf{Актуальность и важность темы.}

То же самое, что и в аннотации, но более подробно.

\textbf{Цель и задачи.}

То же самое.

\textbf{Критерии достижения задач.}

Как вы покажете, что каждая задача была выполнена.

\textbf{Методологическая и технологическая база.}

Какие библиотеки / инструменты / ресурсы вы использовали.

Какой подход вы выбрали.

Возможно, вдохновения.

\textbf{Экспериментальная база исследования.}

Какие данные, сценарии, пользователи, среды, наборы тестов или другие экспериментальные условия использовались.

\textbf{Научная новизна / оригинальность.}

Чем ваша диссертация / продукт отличается от других существующих приложений / решений / исследований.

\textbf{Конкретный вклад автора.}

Какие части результатов являются вашим личным вкладом.

\textbf{Теоретическая значимость полученных результатов.}

Что результаты уточняют, объясняют или систематизируют на теоретическом уровне.

\textbf{Практическая ценность.}

Для чего может быть использовано приложение.
То, что оно помогло вам изучить эти технологии, недостаточно,
нужно указать что-то помимо этого.

\textbf{Апробация, публикация и внедрение результатов.}

Где результаты были представлены, опубликованы, проверены или внедрены.

Можно также сделать такой список:
\begin{itemize}
    \item Ценность A
    \item Ценность B
    \item Ценность C
\end{itemize}

\textbf{Краткое содержание диссертации.}

Первая глава, \nameref{intro_chapter_title}, представляет общую / теоретическую информацию о ...

Вторая глава, \nameref{architecture_chapter_title}, описывает реализацию, ...

Третья глава, \nameref{implementation_chapter_title}, следует за реализацией ...

\chapter{Название теоретической главы}\label{intro_chapter_title}

\section{Синтаксис \LaTeX{}}

\subsection{Введение}

Здесь демонстрируется базовый синтаксис \LaTeX{}.

Очевидно, что подглавы такого размера неприемлемы,
нужно более подробно аргументировать вещи.

Например, можно написать, что базовый синтаксис важен
для написания даже примитивной работы на \LaTeX{}, и поскольку
цель — изучить \LaTeX{} на достаточно продвинутом уровне для написания
диссертации, эта информация необходима для перехода к более глубоким идеям.

\subsection{Сокращения}

\verb|\acronymsChapter| это команда, которая вставляет список сокращений.
Когда у сокращений длинные определения, они автоматичекски красиво отформатируются в таблице,
с минимальным количеством переходов на новую строку.

Сами сокращения определяются до начала документа, используя \verb|\acro| или \verb|\newacronym|.

\verb|\acro| --- это переопределенная команда, использованная в шаблоне.
Первый параметр в квадратных скобках позволяет дать идентификатор сокращению, используя который
будет происходить отсылка к нему в исходном файле работы.
Иначе, имя идентификатора будет то же, что и сокращение.

\verb|\newacronym| ожидает обязательно отдельного идентификатора, что нужно редко, 
поэтому \verb|\acro| использовать проще.

Если определенные сокращения не указать через 
\verb|\gls| или \verb|\acrshort| или аналогичный макрос в работе,
они не появятся в списке сокращений.
По этой причине здесь упоминается \acrshort{LONG TEST}.

\subsection{Цитирование источников}

Цитирование всего источника делается после упоминания, вот так\cite{gif_unusable_reason}.

Можно ссылаться на конкретные страницы или параграфы, вот так\cite[12.2.]{png_spec}.

\subsection{Синтаксис списков}

Ненумерованные. Используются для представления связанной информации:

\begin{itemize}
    \item
        Изображения с палитрой до 256 цветов.
        
    \item
        Возможность потоковой передачи:
        файлы могут читаться и записываться последовательно, что позволяет использовать формат
        PNG как протокол связи для динамической генерации и отображения изображений.

    \item
        Прогрессивное отображение: подготовленный файл изображения может быть
        отображён по мере получения по каналу связи,
        предоставляя очень быстро изображение с низким разрешением,
        с последующим постепенным улучшением деталей.

    \item
        Прозрачность: части изображения могут быть помечены как прозрачные,
        создавая эффект не прямоугольного изображения.

    \item 
        Дополнительная информация: текстовые комментарии и другие данные могут быть
        сохранены внутри файла изображения.

    \item
        Полная независимость от аппаратного обеспечения и платформы.

    \item
        Эффективное сжатие, 100\% без потерь.
\end{itemize}

Нумерованные. Используются, например, для списка последовательных шагов:

\begin{enumerate}
    \item 
        Изображения с палитрой до 256 цветов.
    \item 
        Возможность потоковой передачи:
        файлы могут читаться и записываться последовательно, что позволяет использовать формат
        PNG как протокол связи для динамической генерации и отображения изображений.
    \item 
        Прогрессивное отображение: подготовленный файл изображения может быть
        отображён по мере получения по каналу связи,
        предоставляя очень быстро изображение с низким разрешением,
        с последующим постепенным улучшением деталей.
    \item 
        Прозрачность: части изображения могут быть помечены как прозрачные,
        создавая эффект не прямоугольного изображения.
\end{enumerate}

\subsection{Стилизация}

\textbf{TEXT Bold Face = Жирный}

\textit{TEXT ITalic = Курсив}

\texttt{TEXT TeleType = Моноширинный шрифт для кода}

Можно \textbf{применять} \textit{в предложении}, или \texttt{\textit{\textbf{комбинировать}}}.

\subsection{Кавычки}

Для коротких цитат используйте \verb|\enquote{...}|:
\enquote{это короткая цитата}.

Можно подчеркивать код, например \verb|#include<iostream> int a{~8};|.

Когда дословный текст нужно использовать в хрупком аргументе, например в заголовке,
подписи или таблице, сохраните его через \verb|\SaveVerb| и вставьте через \verb|\UseVerb|.

Также, \verb|\verb| можно использовать не только с \verb!|! для обозначения начала и конца, то есть \verb!\verb|code|!, но можно и, например, \verb|\verb!code!|. 
Используйте тот символ, который не используется в коде, который необходимо вставить как есть.

\subsection{Таблицы}

Когда у вас есть таблица, необходимо дать на неё ссылку.
Название таблицы должно находится \textbf{над} таблицей.
\Объект{pixel_order_table} предоставляет пример вставки таблиц с автоматической нумерацией.

\insertTable[pixel_order_table]{Таблица порядка пикселей\\вторая строка в названии\\третья строка}{%
  \begin{tabular}{c c c c c c c c}
      1 & 6 & 4 & 6 & 2 & 6 & 4 & 6 \\
      7 & 7 & 7 & 7 & 7 & 7 & 7 & 7 \\
      5 & 6 & 5 & 6 & 5 & 6 & 5 & 6 \\
      7 & 7 & 7 & 7 & 7 & 7 & 7 & 7 \\
      3 & 6 & 4 & 6 & 3 & 6 & 4 & 6 \\
      7 & 7 & 7 & 7 & 7 & 7 & 7 & 7 \\
      5 & 6 & 5 & 6 & 5 & 6 & 5 & 6 \\
      7 & 7 & 7 & 7 & 7 & 7 & 7 & 7 \\
  \end{tabular}
}

\SaveVerb{mintedLibraryName}|minted|

\section[Библиотека minted для кода]{Библиотека \UseVerb{mintedLibraryName} для кода}

\subsection{Общая информация}

Библиотека \verb|minted| используется для стилизации кода.
Блоки кода вставляются через вспомогательные команды из шаблона,
чтобы получить нумерацию, подпись и автоматическое размещение.

\subsection{Прямая вставка блока кода}

\Объект{code_direct_example} показывает прямую вставку блока кода.

\begin{code}[code_direct_example]{cpp}{Пример прямой вставки блока кода}
// любой код здесь
int main()
{
    std::cout << "hello";
}
\end{code}

\subsection{Вставка целого файла}

\Объект{../src/sourcefile.zig} вставляет целый файл.

\insertCodeFile{zig}{../src/sourcefile.zig}{Полный исходный файл}

\subsection{Вставка части файла}

\Объект{code_range_example} вставляет только диапазон строк, переданный через опции \verb|minted|.

\insertCodeFile[code_range_example][firstline=2,lastline=5]{zig}{../src/sourcefile.zig}{Диапазон строк из файла}

\subsection{Вставка сегмента из файла}

Эта функциональность реализована благодаря скрипту \texttt{findSegment.py}.

\verb|minted| добавляет подобный функционал (\verb|rangestartafterstring|, \verb|rangestopbeforestring|),
однако не поддерживает отображение номера строки из файла и вставляет пустую строку в начало сегмента.

Язык программирования определяется из расширения файла. 
Для добаления расширений, которые еще не поддерживаются, изменяйте \texttt{findSegment.py}.

\Объект{code_segment_example} вставляет сегмент, выделенный комментариями в исходном файле.

\insertCodeSegment[code_segment_example]{../src/sourcefile.zig}{example}{Отмеченный сегмент из файла}

\Объект{code_segment_autogobble_example} вставляет сегмент с дополнительным отступом.
Опция \verb|autogobble| убирает общий отступ перед отображаемым кодом.

\insertCodeSegment[code_segment_autogobble_example]{../src/sourcefile.zig}{autogobble_example}{Сегмент с автоматически удаленным отступом}

\subsection{Вставка текста на много строк как код}

\Объект{code_text_example} вставляет обычный текст как блок кода.

\begin{code}[code_text_example]{text}{Многострочный текст как код}
Regular text
Multiline
With any characters |||
\end{code}

Чтобы вставить только текст, без номера строки, необходимо добавить \verb|[linenos=false]|.

\Объект{code_text_no_lines_example} скрывает номера строк через опции \verb|minted|.

\begin{code}[code_text_no_lines_example][linenos=false]{text}{Текст без номеров строк}
Regular text
Multiline
With any characters |||
\end{code}

\subsection{Вставка дословного кода в таблицы}

Для коротких фрагментов кода, которые должны появиться в таблице или её подписи,
сохраните дословный текст перед таблицей и вставьте его через \verb|\UseVerb|.

\SaveVerb{tableHttpRoute}|GET /users/{id}?q=a&b%20c|
\SaveVerb{tableJsonConfig}|{"path": "C:\tmp\{id}", "enabled": true}|

\Объект{inline_code_table_example} показывает вставку фрагментов, сохраненных через \verb|\SaveVerb|.

\insertTable[inline_code_table_example]{Фрагменты кода сохранены: \UseVerb{tableHttpRoute}}{%
  \begin{tabularx}{\textwidth}{l X}
    \toprule
    Контекст & Фрагмент \\
    \midrule
    HTTP-маршрут & \UseVerb{tableHttpRoute} \\
    JSON-конфигурация & \UseVerb{tableJsonConfig} \\
    \bottomrule
  \end{tabularx}
}

\section{Изображения}

\subsection{Введение}

Этот шаблон предлагает только одну функцию для изображений, 
которая также автоматически регистрирует его как фигуру.

Картинка будет установлена \textit{не сразу} после упоминания,
а появится автоматически в наиболее близком месте к точке упоминания,
так, чтобы минимизировать пустое место, 
оставляемое, когда картинка не вмещается.

Можно использовать и что-либо другое, если функция вам не подходит.
Для стандартизации, однако, хорошо предложить свой метод через \gls{PR}, чтобы добавить его в шаблон.

\subsection{Центрированное изображение}

Когда вы добавляете изображение, делайте на него ссылку в тексте.

На \объекте{interface_sketch.png} представлен эскиз графического интерфейса.

\insertImage{interface_sketch.png}{Эскиз графического интерфейса}

Этот метод ограничивает максимальную ширину и высоту изображений, 
не давая им выйти за экран и не изменяя их соотношений.
Например, на \объекте{wide.png} широкая картинка не выйдет за пределы,
а на \объекте{long.png} она не уйдет за страницу.

\insertImage{wide.png}{Широкая полоса}

\insertImage{long.png}{Длинная полоса}

\section{Приложения}

\subsection{Выбор между приложением и прямой вставкой кода}

Когда содержание кода не является строго необходимым
для его объяснения на высоком уровне в тексте.

Также, когда код слишком большой.
Попробуйте включать прямо в текст не более 50 строк.

\subsection{Как сделать ссылку}

См. приложение \ref{appendix:example_min}.

В приложении \ref{unnumbered_appendix} представлено приложение без нумерации предметов внутри.

В приложении \ref{numbered_appendix} представлено приложение с нумерацией предметов внутри,
а именно:

\Объекты{cat_1}{cat_2} демонстрируют котов:
\begin{itemize}
  \item \Объект{cat_1} демонстрирует кота номер 1
  \item \Объект{cat_2} демонстрирует кота номер 2
\end{itemize}

В приложении \ref{whole_appendix_reference} представлен пример, где
ссылка на всё приложение покрывает пронумерованные изображения, таблицы и блоки кода внутри него.

\chapterConclusionSection{intro_chapter_title}

Что обсуждалось в этой главе, кратко.

\chapter{Проектирование приложения}\label{architecture_chapter_title}

\section{Введение}

Здесь включите информацию о структуре приложения на высоком уровне,
возможно, некоторые диаграммы, почему вы использовали выбранные технологии и т.д.

\chapterConclusionSection{architecture_chapter_title}

\chapter{Реализация системы}\label{implementation_chapter_title}

Здесь представьте более подробную информацию о конкретных модулях реализованной системы.
Какие проблемы вы встретили и как их преодолели.

Как вы использовали важные фичи библиотек.

Продемонстрируйте ваше приложение, с изображениями — либо непосредственно в тексте,
если они важны для понимания и более специфичны,
либо поместите их в приложение, и вставьте ссылки в текст.

\section{Введение}

Здесь включите информацию о структуре приложения на высоком уровне,
возможно, некоторые диаграммы, почему вы использовали выбранные технологии и т.д.

\chapterConclusionSection{implementation_chapter_title}

В этой главе обсуждалось ...

\finalConclusionChapter

Кратко повторите самое интересное.
Добавьте, что можно улучшить или где продукт может быть использован в будущем.

Пример завершённой диссертации:
\begin{itemize}
  \item Исходный код (программа, текст диссертации на \LaTeX{}): \url{https://github.com/AntonC9018/thesis-png}
  \item PDF: \url{https://drive.google.com/file/d/1ZiGQt6PvUm3FGY8oyIlVsERc4PvZHJu0/view?usp=drive_link}
\end{itemize}

Эти два пункта ниже включайте всегда:

Весь исходный код, включая программный код и текст работы
в форме до рендеринга, доступен на GitHub по следующей ссылке: \github.

Полученные результаты были представлены на \textbf{\conferencesList}\cite{self}.

\bibliographyChapter

\appendixChapter

\section{Функция min в тестовой программе на JavaScript}\label{appendix:example_min}
\insertCodeFile*{js}{../src/appendix_example.js}

\section{Без нумерации}\label{unnumbered_appendix}

\insertImage*{cat_1.jpg}{Кот}

Когда прямая ссылка идет только на все приложение,
а не на индивидуальные картинки, таблицы или блоки кода изнутри,
их не обязательно нумеровать.

\insertTable*{Пример таблицы без нумерации}{%
  \begin{tabular}{c c c}
    Hello & 1 & 2 \\
    A & B & C \\
    D & E & F
  \end{tabular}
}

Для прямого ненумерованного кода в приложениях можно использовать окружение \verb|codeWithoutLabel|.

\begin{codeWithoutLabel}{js}
function min(a, b) {
  return a < b ? a : b;
}
\end{codeWithoutLabel}

\section{Картинка с нумерацией}\label{numbered_appendix}

Это есть смысл делать, если в приложении больше 1 предмета.

\insertImage[cat_1]{cat_1.jpg}{Кот 1}

\insertImage[cat_2]{cat_2.jpg}{Кот 2}

\section{Ссылка на всё приложение}\label{whole_appendix_reference}

В этом случае пронумерованные элементы внутри приложения покрываются ссылкой на приложение.

\insertImage[whole_appendix_cat]{cat_1.jpg}{Кот, покрытый ссылкой на приложение}

\insertTable[whole_appendix_table]{Таблица, покрытая ссылкой на приложение}{%
  \begin{tabular}{c c}
    Элемент & Значение \\
    Изображение & 1 \\
    Таблица & 2
  \end{tabular}
}

\insertCodeFile[whole_appendix_code][]{js}{../src/appendix_example.js}{Код, покрытый ссылкой на приложение}


\ifdefstring{\documentIsPractice}{false}{\declarationPage{}}{}

\end{document}
% vim: fdm=syntax
